% Part: second-order-logic
% Chapter: sol-and-sets
% Section: introduction

\documentclass[../../../include/open-logic-section]{subfiles}

\begin{document}

\olfileid{sol}{set}{int}

\olsection{ভূমিকা}

দ্বিতীয়-ক্রমের যুক্তিবিদ্যায় যেহেতু সংজ্ঞাক্ষেত্রের উপসেটের পাশাপাশি
অপেক্ষকের উপরও পরিমাণায়ন করা যায়, তাই আশা করাই স্বাভাবিক যে অন্তত কিছু
পরিমাণ সেটতত্ত্ব দ্বিতীয়-ক্রমের যুক্তিবিদ্যায় চর্চা করা যাবে। ``চর্চা
করা'' বলতে আমরা বোঝাই, দ্বিতীয়-ক্রমের যুক্তিবিদ্যায় সেটতাত্ত্বিক ধর্ম ও
বক্তব্য প্রকাশ করা সম্ভব, এবং সেটের জন্য কোনো বিশেষ, যুক্তিবহির্ভূত শব্দভাণ্ডার
(যেমন সেটতত্ত্বের সদস্যতা-!!{predicate}) ছাড়াই তা সম্ভব। যেমন, আমরা সেটের
সংযোগ ও ছেদ এবং উপসেট-সম্পর্ক সংজ্ঞায়িত করতে পারি; আবার সেটের আকারও তুলনা
করতে পারি এবং কান্টরের উপপাদ্যের মতো ফলও বলতে পারি।

\end{document}
